\documentclass[11pt,a4paper]{article}
\usepackage[utf8]{inputenc}
\usepackage[brazil]{babel}
\usepackage{graphicx}
\usepackage{amssymb, amsfonts, amsmath}
\usepackage{color}
\usepackage{float}
\usepackage{enumerate}
%\usepackage{subfigure}
\usepackage[top=2.5cm, bottom=2.5cm, left=1.25cm, right=1.25cm]{geometry}

\newcommand{\sen}{\,\textrm{sen}\,}
\newcommand{\tg}{\,\textrm{tg}\,}

\begin{document}
\pagestyle{empty}

\begin{center}
\begin{tabular}{ccc}
\begin{tabular}{c}
\includegraphics[scale=0.25]{../../biblioteca/imagem/brasao-de-armas-brasil} \\
\end{tabular} & 
\begin{tabular}{c}
Ministério da Educação \\
Universidade Federal dos Vales do Jequitinhonha e Mucuri \\
Faculdade de Ciências Sociais, Aplicadas e Exatas - FACSAE \\
Departamento de Ciências Exatas - DCEX \\
Disciplina: Cálculo Numérico\\
Prof.: Luiz C. M. de Aquino\\
\end{tabular} &
\begin{tabular}{c}
\includegraphics[scale=0.25]{../../biblioteca/imagem/logo-ufvjm} \\
\end{tabular}
\end{tabular}
\end{center}

\begin{center}
 \textbf{Lista de Exercícios IV}
\end{center}

\textbf{Observação:} Nos exercícios 1 e 2 use o Método de Eliminação Gaussiana para resolver o sistema de equações.

\begin{enumerate}
  \item Um fabricante de móveis produz cadeiras, mesinhas de centro e mesas
de jantar. Cada cadeira leva 10 minutos para ser lixada, 6 minutos
para ser tingida e 12 minutos para ser envernizada. Cada mesinha de
centro leva 12 minutos para ser lixada, 8 minutos para ser tingida
e 12 minutos para ser envernizada. Cada mesa de jantar leva 15 minutos
para ser lixada, 12 minutos para ser tingida e 18 minutos para ser
envernizada. A bancada para lixar fica disponível 1.340 minutos por
semana, a bancada para tingir 940 minutos por semana e a bancada para
envernizar 1.560 minutos por semana. Quantos móveis devem ser fabricados
(por semana) de cada tipo para que as bancadas sejam plenamente utilizadas?

 \item Considere o problema do circuito hidráulico mostrado na Figura~\ref{fig:tub}. Este sistema está alimentado por um reservatório 
cuja a pressão é mantida constante e igual a $P_r = 10$. As saídas das tubulações desembocam na atmosfera, onde a pressão é 
considerada nula (isto é, $P_a = 0$). Deste modo, a vazão $Q_i$ da $i$-ésima tubulação depende da diferença de pressão $\Delta P_i$ 
de tal modo que $$Q_i = K_iL_i\Delta P_i,$$ onde $K_i$ é a resistência hidráulica e $L_i$ o comprimento da tubulação. Por exemplo, para a 
tubulação $8$ temos que $Q_8 = K_8L_8\Delta P_8$, sendo que $\Delta P_8 = P_1 - P_4$ (ou seja, a pressão que ``entra'' na tubulação pela 
bifurcação $1$ menos a pressão que ``sai'' da tubulação pela bifurcação $4$). Por outro lado, sabe-se que em cada bifurcação a soma das vazões 
deve ser nula. Por exemplo, na bifurcação $4$ temos que $Q_8 - Q_6 - Q_7 = 0$ (aqui note que a vazão que ``entra'' na bifurcação é considerada positiva, 
enquanto que a que ``sai'' é considerada negativa). Considerando essas informações e os dados da Tabela~\ref{tab:dados}, determine as vazões em cada tubulação e as pressões 
em cada bifurcação.

\begin{figure}[H]
 \centering
 \includegraphics[scale=0.75]{imagem/Tubulacao.png}
 \caption{Esquema do circuito hidráulico.}
 \label{fig:tub}
\end{figure}

\begin{table}[H]
 \centering
 \begin{tabular}{c|c|c}
  Tubulação $i$ & $K_i$ & $L_i$ \\ \hline
  1 & 0,02 & 1,0 \\ \hline
  2 & 0,005 & 2,0 \\ \hline
  3 & 0,085 & 0,5 \\ \hline
  4 & 0,02 & 1,0 \\ \hline
  5 & 0,075 & 0,5 \\ \hline
  6 & 0,085 & 0,5 \\ \hline
  7 & 0,015 & 2,0 \\ \hline
  8 & 0,01 & 1,0 
 \end{tabular}
 \caption{Resistência hidráulica e comprimento das tubulações.}
 \label{tab:dados}
\end{table}
 
  \item Pesquise uma maneira de adaptar o Método de Eliminação Gaussiana para calcular o determinante de uma matriz. Em seguida, use este método para calcular 
o determinante da matriz:
$$A=\begin{bmatrix}3 & - 5 & 1 \\ -1 & 1 & 3 \\ -7 & 3 & -6\end{bmatrix}.$$

  \item Considere as matrizes:
$$A=\begin{bmatrix}1 & - 1 & 2 \\ 3 & 7 & -4 \\ 5 & -2 & 1\end{bmatrix}, \, 
x=\begin{bmatrix}x_1 \\ x_2 \\ x_3\end{bmatrix}, \, 
b=\begin{bmatrix}-2 \\ 1 \\ 5\end{bmatrix} \textrm{ e }\,
c=\begin{bmatrix}1 \\ -1 \\ 3\end{bmatrix}. 
$$

  \begin{enumerate}
    \item Determine a fatoração $LU$ de $A$.
    \item Use a fatoração do item anterior para resolver os sistemas $Ax = b$ e $Ax = c$.
  \end{enumerate}

  \item Suponha que uma matriz $A$ foi fatorada no formato $LU$ (sem utilizar pivoteamento). Preencha os espaços em branco abaixo de modo a determinar 
as matrizes $A$, $L$ e $U$.

$$
A=\begin{bmatrix}\square & \square & \square \\ 1 & 0 & \square \\ \square & -4 & -3\end{bmatrix};\quad
L=\begin{bmatrix}1 & 0 & 0 \\ \square & 1 & 0 \\ -5 & \square & 1\end{bmatrix};\quad
U=\begin{bmatrix}-2 & \square & 4 \\ 0 & \square & 3 \\ 0 & 0 & -1\end{bmatrix}
$$

\end{enumerate}

\begin{center}
\textbf{Gabarito}
\end{center} 
\textbf{[1]} $50$ cadeiras, $20$ mesinhas de centro e $40$ mesas de jantar. 
\textbf{[2]} $P_1 = 5,47246730738932$, $P_2 = 0,919331967858831$, $P_3 = 0,596344729793603$, $P_4 = 0,970537261698440$. 
$Q_1 = 0,0905506538522137$, $Q_2 = 0,0455313533953049$, $Q_3 = 0,0390716086340003$, $Q_4 = 0,00645974476130455$, 
$Q_5 = 0,0223629273672601$, \\ $Q_6 = 0,0159031826059556$, $Q_7 = 0,0291161178509532$, $Q_8 = 0,0450193004569088$.
\textbf{[3]} $\det A = 94$.
\textbf{[4]} (a) $L = \begin{bmatrix}1 & 0 & 0 \\ 3 & 1 & 0 \\ 5 & \frac{3}{10} & 1\end{bmatrix}$ e 
$U = \begin{bmatrix} 1 & -1 & 2 \\ 0 & 10 & -10 \\ 0 & 0 & -6 \end{bmatrix}$. 
(b) $x = 
\begin{bmatrix}
\frac{17}{20} \\
-\frac{29}{20} \\
-\frac{43}{20}
\end{bmatrix}
$. 
$x = 
\begin{bmatrix}
\frac{7}{15} \\
-\frac{4}{15} \\
\frac{2}{15}
\end{bmatrix}
$. 
\textbf{[5]} $A = \begin{bmatrix}-2 & 2 & 4 \\ 1 & 0 & 1\\ 10 & -4 & -3\end{bmatrix}$, 
$L = \begin{bmatrix}1 & 0 & 0 \\ -\frac{1}{2} & 1 & 0 \\ -5 & 6 & 1\end{bmatrix}$ e  
$U = \begin{bmatrix}-2 & 2 & 4 \\ 0 & 1 & 3 \\ 0 & 0 & -1\end{bmatrix}$
\end{document}
